% Generated draft: requires mathematical and source review.
\documentclass[11pt]{article}
\usepackage[margin=25mm]{geometry}
\usepackage{fontspec}
\setmainfont{Latin Modern Roman}
\usepackage{amsmath,amssymb,mathtools,unicode-math}
\setmathfont{Latin Modern Math}
\usepackage{enumitem,needspace,xurl,hyperref}
\setlength{\parindent}{0pt}
\setlength{\parskip}{5pt}
\setlength{\emergencystretch}{3em}
\newcommand{\N}{\mathbb{N}}
\newcommand{\Z}{\mathbb{Z}}
\newcommand{\Q}{\mathbb{Q}}
\newcommand{\R}{\mathbb{R}}
\newcommand{\C}{\mathbb{C}}
\newcommand{\F}{\mathbb{F}}
\newcommand{\A}{\mathbb{A}}
\newcommand{\PP}{\mathbb{P}}
\newcommand{\E}{\mathbb{E}}
\newcommand{\eps}{\varepsilon}
\newcommand{\dd}{\,\mathrm{d}}
\begin{document}
\section*{UnsolvedMath ID: 2303005}
\textbf{Source problem code:} catalogue entry
\subsection{Definitions and mathematical statement}
Let $\mathbb D=\{z\in\C:|z|<1\}$, and let $u>0$ be subharmonic on $\mathbb D$ (upper semicontinuous and satisfying the sub-mean inequality). The source assumes a sequence of arcs $\gamma_n$ tending to the boundary arc $\{e^{i\theta}:\alpha\leq\theta\leq\beta\}$, with a common upper bound $u\leq M$ on them. Its printed index range is only ``$n=1,2$,'' although a sequence tending to the boundary suggests that all $n\geq1$ may be intended. The topology and geometric conditions on these arcs are not supplied.

The source next prints
\[
\int_0^1(1-r)u(re^{i\theta})\,d\theta<\infty
\]
for $\theta$ in a set $E$ dense in $(\alpha,\beta)$. As written, the differential and free parameter are inconsistent. A plausible radial condition is
\[
\int_0^1(1-r)u(re^{i\theta})\,dr<\infty
\quad\text{for each }\theta\in E,
\]
but replacing $d\theta$ by $dr$ requires verification and is not treated as a confirmed transcription.

Under the intended hypotheses, the source attributes to Maclane the conclusion that for each $0<\delta<(\beta-\alpha)/2$ there is $C_\delta<\infty$ with
\[
u(re^{i\theta})\leq C_\delta
\quad\text{whenever }0\leq r<1,\quad
\alpha+\delta\leq\theta\leq\beta-\delta.
\]
It says this conclusion holds in particular under
\[
\int_0^1(1-r)B(r,u)\,dr<\infty.
\]
The conventional meaning of $B(r,u)$ would be $\max_{|z|=r}u(z)$, but its definition is not reproduced in this record.

Problem 3.5 asks whether the growth conditions referenced as source labels 3.3 and 3.4 can be weakened while retaining the same uniform bound on every closed angular subinterval. The labels are unresolved, so their exact identification with the displayed hypotheses must be checked. The question does not propose a specific weaker weight, integrability class, or angular condition. The supplied summary reports it open as of the 2018 edition, with no recent resolution located.
\subsection{Short English statement}
The problem asks whether growth assumptions guaranteeing uniform bounds for a positive subharmonic function near an interior portion of a boundary arc can be weakened. The extracted radial integral has the wrong-looking differential, the arc index range is incomplete, and the referenced condition labels are unresolved.
\subsection{Sources}
\begin{itemize}
\item Local UnsolvedMath registry, permanent ID 2303005, exact source code catalogue entry, ``Research Problems in Function Theory --- Problem 3.5''; status open, classification OPEN-TRIAGE.
\item Hayman, \emph{Research Problems in Function Theory} (2018), Problem 3.5, \url{https://arxiv.org/abs/1809.07200}. The record identifies extraction from public arXiv source TeX, accessed July 29, 2026; the 2018 update was retained for status triage.
\item The statement attributes the sufficient-condition theorem to Maclane but supplies no full citation. Source labels 3.3 and 3.4 are unresolved in the extraction.
\item The supplied summary reports openness as of the 2018 edition and no recent resolution located, with no additional authoritative citation or precise search date.
\end{itemize}
\end{document}
